\chapter*{Introduction}
\addcontentsline{toc}{chapter}{Introduction}

\label{ch:introduction}
Heat exchangers are crucial in industrial and energy systems, where precise temperature control is often needed. Controlling them with conventional industrial methods can be challenging due to their complex behavior, operational limitations, and poor response times. As a result, more advanced control strategies are being used.

Model Predictive Control (MPC) is a widely adopted method in process engineering, providing a systematic way to control constrained Single-Input Single-Output (SISO) and Multiple-Input Multiple-Output (MIMO) systems through finite-horizon optimization that explicitly accounts for system constraints, making it ideal for various industrial applications such as thermal processes \cite{rawlings_mpc}.

MPC has also been reported in heat-exchanger control applications, where the main goal is usually outlet temperature regulation under input and output constraints. These studies show that predictive control can improve temperature regulation compared with simpler feedback strategies, especially when the operation is close to actuator or safety limits~\cite{mpc_heat_exchanger, grune2016nonlinear}.


Classical linear models continue to be widely used in industrial settings because they are simple and straightforward to identify. For thermal systems like heat exchangers, step-response tests are often performed to develop low-order linear models. Such models generally capture the main dynamics with a steady-state gain and a single time constant, which is usually adequate for control within a limited operating range \cite{morari1999}. In SISO heat exchanger applications, low-order linear models are especially common \cite{mpc_heat_exchanger}.


These models perform well around a nominal operating point, but when nonlinear factors are introduced, their accuracy decreases. Because of this constraint, more sophisticated modeling techniques for predictive control have been used \cite{morari1999}, and research interest in data-driven and nonlinear modeling techniques has grown.

Data-driven modeling offers an alternative to traditional first-principles and step-response-based methods. Instead of relying solely on physical laws or simplified assumptions, data-driven techniques build system models directly from measured input–output data. This can be especially valuable for thermal systems, where nonlinear heat transfer, changing operating conditions, and unmeasured disturbances make purely physical modeling difficult \cite{morari1999}.


Several recent works show the use of data-driven models. A study investigated data-driven modeling of heat pumps and thermal storage units for MPC and compared static models with dynamic state-space approaches based on Dynamic Mode Decomposition with control (DMDc) and Sparse Identification of Nonlinear Dynamics with control (SINDyc). In that work, the models were evaluated for heating and cooling case studies, including their prediction performance and model complexity. The study further observed that while neural network models can capture nonlinear dynamics, they typically demand more training data and offer less interpretability than simpler state-space methods \cite{brandes2022}.

Another example is a work where support vector regression (SVR) was used for data-driven prediction of steady-state conditions in heat exchangers. While SVR is derived from the support vector machine (SVM) principle, it is specifically designed for regression tasks rather than classification. In this approach, the model learns the relation between measured operating variables and the steady-state output without requiring a detailed analytical description of the heat exchanger. The authors demonstrated that the SVR-based approach could estimate steady-state operating conditions with mean errors under 3.5\%, while also remaining robust to variations in input variables \cite{wu2025svr}.



The Koopman operator framework is one of the methods that provides a theoretical foundation for using linear operators to represent nonlinear dynamical systems within the lifted state space. Koopman \cite{koopman1931} first proposed the idea, and Mezi\'c later developed it and showed how it might be used to analyze the behavior of nonlinear systems \cite{mezic2005}.

An essential feature of the Koopman framework is the lifting function $g(\cdot)$, which transforms the original system states or measurements into a higher-dimensional space. In this transformed space, nonlinear dynamics can be modeled using a linear prediction approach. Consequently, choosing or designing an appropriate lifting function is vital for achieving an accurate Koopman model.


Various strategies for constructing the lifting function $g(\cdot)$ have been proposed in the literature. Predefined dictionaries of basis functions, including polynomials, were used in early data-driven methods. Extended Dynamic Mode Decomposition (EDMD) is a well-known example, approximating the Koopman operator using linear regression in a specified function space \cite{williams2015}. The effectiveness of such methods depends on the selected dictionary and can be constrained for complex or high-dimensional systems.

Recent studies have increasingly emphasized learning the lifting function directly from data. It has been shown that sparse data-driven representations can identify system dynamics in the absence of explicit physical models \cite{brunton2016}. Building on this foundation, Koopman-based linear predictors have been developed for nonlinear systems, and their applicability to optimal control problems has been established \cite{korda2018}. Neural-network-based lifting functions have further advanced this approach by providing flexible representations of nonlinear dynamics while retaining a linear structure suitable for control design \cite{koopman_mpc_review}.


The Koopman framework has also been incorporated into MPC frameworks \cite{korda2018}.
In Koopman-based MPC, the learned linear model is employed directly within the MPC optimization problem, allowing nonlinear system behaviors to be represented using a linear prediction model. 
This integration enables direct comparison of Koopman-based MPC and classical first-order linear MPC under identical constraints and cost functions. This thesis investigates how Koopman-based modeling influences control performance.


Other recent studies have shown how Koopman-based MPC may be used in thermal systems. For example, Koopman–Model Predictive Control was implemented in a warm-water supply system, where a data-driven linear Koopman model regulated outlet temperature under variable demand and operating conditions, resulting in improved tracking and constraint handling compared to conventional control approaches \cite{miyashita2022koopman}.
Other studies have explored vehicle thermal management, utilizing Koopman-based MPC to regulate temperatures in automotive thermal subsystems with nonlinear heat transfer characteristics, all while ensuring actuator and safety constraints are satisfied \cite{chen2025koopman_tms}.





